\section{Additional Related Work}

\textbf{Brain Rot Effects.}
Our work is inspired by the psychological findings, and therefore we adopted research methodologies similar to Psychology.
A study on 15-16 year olds published in the Journal of American Medical Association found “significant association between higher frequency of modern digital media use and subsequent symptoms of ADHD” \citep{ra2018association}.
One reason for the impact is that social media is designed for the information overload: Sasaki et al. found that the more the number of Twitter friends you have, the higher the risk of information overload~\citep{sasaki2015anatomy}. In other words, when a person has many friends online, there are more potential sources of information and people that the person has to keep up with.
% \textbf{Memory.} The online world often act as a form of external memory or “transactive memory”
Media multitasking is also associated with distractibility and increased prefrontal activity in adolescents and young adults \citep{moisala2016media}.
Recently, Brain Rot was connected with GenAI. \cite{eliot2024brainrot} points out that the prevalence of GenAI increases the risks of human brain rot, though GenAI could also be useful for reducing Brain Rot.
Yet, there is no study on whether LLMs can get brain rot like humans. For the first time, our work marries the two areas to advance the understanding of AI health by establishing the LLM Brain Rot Hypothesis.
